\chapter{Classification and Coverage}
\label{app:classification}

\section{Evidence labels}

Every substantive claim in this documentation carries one of five labels. The
labels are used literally as \textbf{Status:} markers throughout.

\begin{longtable}{p{3.4cm}p{11cm}}
\caption{Evidence labels}\\
\toprule
\textbf{Label} & \textbf{Meaning} \\
\midrule

IMPLEMENTED & The behaviour exists in the repository and the claim is supported by
a specific file, configuration or artefact. \\
INFERRED & A conclusion drawn from structure or behaviour that is not stated
explicitly. The reasoning is given. \\
RECOMMENDED & A proposed improvement that is \emph{not} present. Never presented
as if it existed. \\
UNKNOWN & The claim cannot be determined from the repository (for example
external hosting state). \\
NOT APPLICABLE & The concept does not exist in this product. Never used to hide a
gap --- gaps are labelled as such. \\
\bottomrule
\end{longtable}

Supporting conventions: security findings use \emph{Finding / Severity / Evidence
/ Impact / Recommendation}; maintenance items use \emph{Open / Accepted / Pass}
or \emph{Open / Unknown}; maintenance registers use identifiers (SEC, KI, TD,
PERF, SCALE, SEO, AEO, GEO, RM, RN).

\section{Coverage of requested document types}

The table below records every document type that was requested and states how it
was handled. Section types marked \textbf{NOT APPLICABLE} were deliberately
\emph{not} written as standalone chapters, because writing them would mean
padding a chapter with no repository evidence. Their non-applicability is stated
explicitly so nothing is silently omitted.

\begin{longtable}{p{4.6cm}p{2.7cm}p{7.1cm}}
\caption{Requested document types and their handling}\\
\toprule
\textbf{Requested type} & \textbf{Handling} & \textbf{Location / reason} \\
\midrule

Project overview & Written & Chapter~\ref{ch:project-overview} \\
Project description & Written & Chapter~\ref{ch:project-description} \\
Executive summary & Written & Chapter~\ref{ch:executive-summary} \\
Business requirements (BRD) & Written & Chapter~\ref{ch:brd} \\
Product requirements (PRD) & Written & Chapter~\ref{ch:prd} \\
Functional requirements & Written & Chapter~\ref{ch:functional-requirements} \\
Non-functional requirements & Written & Chapter~\ref{ch:non-functional-requirements} \\
Technical requirements & Written & Chapter~\ref{ch:technical-requirements} \\
Software architecture & Written & Chapter~\ref{ch:software-architecture} \\
System architecture & Written & Chapter~\ref{ch:system-architecture} \\
High-level design & Written & Chapter~\ref{ch:high-level-design} \\
Low-level design & Written & Chapter~\ref{ch:low-level-design} \\
Frontend architecture & Written & Chapter~\ref{ch:frontend-architecture} \\
Backend architecture & NOT APPLICABLE & No backend or server application exists
in the repository; the only server is a local static JSON fixture
(\texttt{scripts/serve.mjs}). \\
Database design & NOT APPLICABLE & There is no database, ORM or persistence layer.
The only persisted state is a JSON manifest in the user's project. \\
API documentation & Written & Chapter~\ref{ch:api-documentation} (registry HTTP
contract and the core TypeScript API). \\
Authentication & Written & Chapter~\ref{ch:authentication} (no real auth in the
product; the \emph{generated} auth is analysed). \\
Authorization & Written & Chapter~\ref{ch:authorization} (largely not applicable;
records the absence of roles/permissions). \\
User flows & Written & Chapter~\ref{ch:user-flows} \\
Data flows & Written & Chapter~\ref{ch:data-flows} \\
UI/UX & Written & Chapter~\ref{ch:ui-ux} \\
Design system & Written & Chapter~\ref{ch:design-system} \\
Security & Written & Chapter~\ref{ch:security} \\
Performance & Written & Chapter~\ref{ch:performance} \\
Scalability & Written & Chapter~\ref{ch:scalability} \\
Testing & Written & Chapter~\ref{ch:testing} \\
Deployment & Written & Chapter~\ref{ch:deployment} \\
CI/CD & Written & Chapter~\ref{ch:cicd} \\
Integrations & Written & Chapter~\ref{ch:integrations} \\
Configuration & Written & Chapter~\ref{ch:configuration} \\
SEO & Written & Chapter~\ref{ch:seo} \\
AEO & Written & Chapter~\ref{ch:aeo} \\
GEO & Written & Chapter~\ref{ch:geo} \\
Known issues & Written & Chapter~\ref{ch:known-issues} \\
Technical debt & Written & Chapter~\ref{ch:technical-debt} \\
Roadmap & Written & Chapter~\ref{ch:roadmap} \\
Development guidelines & Written & Chapter~\ref{ch:development-guidelines} \\
Contribution guide & Written & Chapter~\ref{ch:contribution-guide} \\
Glossary & Written & Chapter~\ref{ch:glossary} \\
\bottomrule
\end{longtable}

\section{Additional chapters written}

Two additional chapters were written because the repository contains material
that the requested list did not cover explicitly and that would otherwise go
undocumented:

\begin{itemize}
  \item \textbf{CLI reference} (Chapter~\ref{ch:cli-reference}) --- the CLI is a
        primary product surface with five commands and a precise option set.
  \item \textbf{Registry and recipes}
        (Chapter~\ref{ch:registry-and-recipes}) --- the recipe schema and catalog
        are the product's core data model.
\end{itemize}

Three appendices were added: the catalog reference
(Chapter~\ref{app:catalog}), the environment and toolchain
(Chapter~\ref{app:environment}), and this classification appendix. A file index
follows (Chapter~\ref{app:file-index}).

\section{Applicability summary}

\begin{longtable}{p{4.4cm}p{2.4cm}p{7.6cm}}
\caption{Areas that are not applicable to this product}\\
\toprule
\textbf{Area} & \textbf{Verdict} & \textbf{Basis} \\
\midrule

Backend / server application & NOT APPLICABLE & No server code. \\
Database / migrations & NOT APPLICABLE & No database. \\
User accounts / tenancy & NOT APPLICABLE & No accounts. \\
Role-based access control & NOT APPLICABLE & No roles or permissions. \\
Web / browser UI & NOT APPLICABLE & No web application. \\
Public website / landing page & NOT APPLICABLE & No site; registry is JSON. \\
Telemetry / analytics & NOT APPLICABLE & None present (and none is a gap to fill). \\
Payment / billing & NOT APPLICABLE & Product is free/open-source packages plus a
privately licensed extension; no transactions in code. \\
Containerisation / IaC & NOT APPLICABLE & Static hosting only. \\
Real-time / streaming & NOT APPLICABLE & No such feature. \\
\bottomrule
\end{longtable}

\section{Reading guide}

\begin{itemize}
  \item For a business reader: Chapters~\ref{ch:project-overview}--\ref{ch:prd},
        plus this appendix and Chapter~\ref{ch:known-issues}.
  \item For an engineer integrating or extending: Chapters
        \ref{ch:software-architecture}--\ref{ch:frontend-architecture},
        \ref{ch:cli-reference}--\ref{ch:data-flows}, and
        Chapters~\ref{ch:development-guidelines}--\ref{ch:contribution-guide}.
  \item For an assessor: Chapters~\ref{ch:security}, \ref{ch:performance},
        \ref{ch:scalability}, \ref{ch:testing}, \ref{ch:cicd}, and
        Chapters~\ref{ch:known-issues}--\ref{ch:roadmap}.
\end{itemize}

\textbf{Status:} IMPLEMENTED (this appendix describes the documentation itself and
the applicability of each area, all verifiable from the repository).
